% !TeX root = main.tex
% Deckblattdaten in protokolldaten.tex bearbeiten.
\documentclass[11pt,a4paper]{scrartcl}
\input{vorlage/praeambel}
% Zentrale Angaben für das Deckblatt. Sonderzeichen: \& statt &, \_ statt _.
\newcommand{\Hochschule}{Technische Universität Wien}
\newcommand{\Institut}{Institut für Mechatronik und Leistungselektronik}
\newcommand{\Lehrveranstaltung}{369.011 Messtechnik Labor}
\newcommand{\Semester}{Semester 5}
\newcommand{\Versuchsnummer}{Laborübung 3}
\newcommand{\Versuchstitel}{Messbrücken und Messverstärker}
\newcommand{\Versuchsuntertitel}{Untertitel oder kurze Beschreibung}
\newcommand{\Verfasser}{Mader Moritz}
\newcommand{\Matrikelnummer}{12409760}
\newcommand{\Laborgruppe}{6}
\newcommand{\Messplatz}{CA0410-5}
\newcommand{\Laborpartner}{Hu Andy, Plattner Magdalena}
\newcommand{\Versuchsdatum}{10.11.2026}
\newcommand{\Abgabedatum}{Abgabedatum}
% Optional: für ein Deckblatt ohne Logo den Inhalt der Klammern leeren.
\newcommand{\Logodatei}{bilder/mpei.pdf}

\KOMAoptions{listof=totoc,bibliography=totoc}
\hypersetup{pdftitle={Laborprotokoll \Versuchsnummer: \Versuchstitel},
  pdfauthor={\Verfasser},pdfsubject={\Lehrveranstaltung}}

\DeclareNewTOC[
  type=equation,
  listname={Gleichungsverzeichnis}
]{loe}
\setuptoc{loe}{totoc}

\begin{document}
\hypersetup{pageanchor=false}
% !TeX root = ../main.tex
\begin{titlepage}
  \thispagestyle{empty}
  {\sffamily
  \setlength{\parfillskip}{0pt plus 1fil}
  \noindent\begin{minipage}[t]{0.72\textwidth}
    \vspace{0pt}
    {\large\bfseries \Hochschule\par}
    \smallskip
    {\color{Hinweisgrau}\Institut\par}
  \end{minipage}\hfill
  \begin{minipage}[t]{0.24\textwidth}
    \vspace{0pt}
    \raggedleft
    \ifdefempty{\Logodatei}{}{%
      \IfFileExists{\Logodatei}{%
        \includegraphics[width=\linewidth,height=20mm,keepaspectratio]{\Logodatei}%
      }{\textcolor{Hinweisgrau}{\small Logo fehlt}}
    }
  \end{minipage}\par
  \vspace{23mm}
  {\noindent\color{Akzent}\rule{\textwidth}{1.4pt}\par}
  \vspace{7mm}
  {\large\color{Akzent}\MakeUppercase{Laborprotokoll}\par}
  \vspace{5mm}
  {\Huge\bfseries \Versuchsnummer\par}
  \vspace{4mm}
  {\LARGE\bfseries \Versuchstitel\par}
  \vspace{4mm}
  %{\large\color{Hinweisgrau}\Versuchsuntertitel\par}
  %\vspace{9mm}
  {\large\Lehrveranstaltung\par}
  %{\color{Hinweisgrau}\Semester\par}
  \vfill
  \renewcommand{\arraystretch}{1.55}
  \begin{tabularx}{\textwidth}{@{}p{43mm}Y@{}}
    \toprule
    Verfasst von & \Verfasser \\
    Matrikelnummer & \Matrikelnummer \\
    Laborgruppe & \Laborgruppe \\
    Laborpartner/innen & \Laborpartner \\
    \midrule
    Versuchsdatum & \Versuchsdatum \\
    Abgabedatum & \Abgabedatum \\
    \bottomrule
  \end{tabularx}
  \vspace{8mm}
  }
\end{titlepage}

% Protokollinhalt direkt in dieser Hauptdatei bearbeiten.
\pagenumbering{Roman}
\hypersetup{pageanchor=true}
\tableofcontents
\clearpage
% Abbildungen mit \caption{...} werden hier automatisch aufgenommen.
\listoffigures
\clearpage

\listofequations
\clearpage

\pagenumbering{arabic}
\pagestyle{scrheadings}
\section{Formalkriterien}
\textbf{Verwendete Dokumentvorlage:}
Dieses Protokoll verwendet eine angepasste Version der LaTeX-Vorlage
\href{https://github.com/bendemaj/Laborprotokoll}{\texttt{bendemaj/Laborprotokoll}}
von GitHub~\cite{vorlage-laborprotokoll}.

\subsection{Titel und Bezeichnung der Teilübung}
\subsection{Datum, Teilnehmer und Messplatz}
\begingroup\setlength{\parfillskip}{0pt}\noindent
\begin{tabularx}{\textwidth}{@{}>{\raggedright\arraybackslash}p{37mm}Y@{}}
Versuchsdatum & \Versuchsdatum \\
Teilnehmer & \Verfasser; \Laborpartner \\
Laborgruppe & \Laborgruppe \\
Laborplatz & \Messplatz \\
\end{tabularx}\par\endgroup
% Foto labor_platz.jpg: CA0410-5 ist auf einer Anschlussleiste sichtbar.
% Der Messplatz wird aus der bereits ausgefuellten protokolldaten.tex uebernommen.

\subsection{Verwendete Geräte}
\label{sec:geraeteliste}
Die verwendeten Geräte und Hilfsmittel sind in \cref{tab:geraete} aufgelistet.
Die Gerätebezeichnungen wurden aus den Fotos übernommen. Fehlende Angaben sind in der Tabelle markiert und werden noch ergänzt.

\begin{table}[htbp]
  \centering
  \caption{Geräte und Hilfsmittel für Laborübung 3.}
  \label{tab:geraete}
  \begingroup
  \small
  \setlength{\tabcolsep}{3pt}
  \renewcommand{\arraystretch}{1.3}
  \begin{tabularx}{\textwidth}{@{}>{\raggedright\arraybackslash}p{30mm}>{\raggedright\arraybackslash}p{46mm}>{\raggedright\arraybackslash}p{29mm}Y@{}}
    \toprule
    Gerät / Hilfsmittel & Hersteller / Typ & Inventarnr. / Kennung & Messbereich / Bemerkung \\
    \midrule
    Digitalmultimeter & Keysight U1232A & \platzhalter{ergänzen} & Diagonalspannung; \platzhalter{Messbereich} \\
    Digitalmultimeter & Neumann 9140 & \platzhalter{ergänzen} & \platzhalter{Messbereiche} \\
    Oszilloskop & Agilent InfiniiVision DSO-X 2002A & CD0408-6 & \platzhalter{Tastköpfe} \\
    Funktionsgenerator & Siglent SDG1025 & CA0406-5 & \platzhalter{Betriebsart} \\
    DC-Netzgerät & RS PRO RS-D3305P & \platzhalter{ergänzen} & \platzhalter{Kanäle} \\
    Widerstandsdekade & Goerz Electro & Nr.~A~2475 & Aufdruck: 529 MB 65/32 \\
    Kapazitätsdekade & Danbridge DK4S & \platzhalter{ergänzen} & \platzhalter{verwendeter Bereich} \\
    Steckbrett & \platzhalter{Typ ergänzen} & CA0406-6 & Zusatzwiderstände und Verdrahtung \\
    Potentiometer & \platzhalter{Typ ergänzen} & \platzhalter{ergänzen} & \platzhalter{Nennwert} \\
    DMS-Wägeaufbau & Laboraufbau & \platzhalter{ergänzen} & Mit Temperatursensoren und Heizung \\
    Messverstärker-Platine & \platzhalter{Bezeichnung} & \platzhalter{ergänzen} & Subtrahier- und Instrumentenverstärker \\
    Weitere Geräte / Zubehör & \platzhalter{ergänzen} & \platzhalter{ergänzen} & \platzhalter{z. B. Referenzgewichte} \\
    \bottomrule
  \end{tabularx}
  \endgroup
\end{table}
% Fotobelege: bilder/used_devices/. Die Messverstaerker-Platine ist in den
% hier bereitgestellten Geraetefotos nicht eindeutig separat dokumentiert.
\subsection{Kurze einführende Beschreibung (optional)}
\subsection{Bestätigung der selbst durchgeführten Messungen und der daraus gewonnenen Ergebnisse}
\subsection{Unterschriften aller Teilnehmer}

\clearpage
\section{DC-Messbrücke}
Ziel dieses Versuchs ist die Bestimmung des parasitären Parallelwiderstands eines Kondensators mit einer DC-Wheatstone-Abgleichmessbrücke.
Der Parallelwiderstand beschreibt die Leckströme durch das nicht ideal isolierende Dielektrikum.

\subsection{Messschaltung}
\Cref{fig:wien_bridge} zeigt eine Kapazitätsmessbrücke nach Wien mit zwei Parallel-RC-Zweigen und zwei rein ohmschen Brückenarmen.
Dabei sind $R_1$ und $C_1$ die einstellbaren Vergleichsgrößen; $R_2$ und $C_2$ beschreiben den zu untersuchenden Kondensator.
Für die DC-Messung wird der Vergleichskondensator $C_1$ weggelassen.
Nach dem Einschwingen fließt durch den idealen kapazitiven Anteil von $C_2$ kein Gleichstrom; der Parallelwiderstand $R_2$ bleibt wirksam.

\begin{figure}[htbp]
  \centering
  \includegraphics[width=0.5\textwidth]{shematics/wien_bridge.png}
  \caption{Kapazitätsmessbrücke nach Wien; für die DC-Messung wird $C_1$ weggelassen.}
  \label{fig:wien_bridge}
\end{figure}

Zum Einstellen des Widerstands $R_1$ stand eine Widerstandsdekade zur Verfügung.
% Foto der verwendeten Widerstandsdekade hier ergänzen.

\subsection{Wahl der Versorgungsspannung}
Zur Bestimmung des Parallelwiderstands $R_2$ wurde eine Gleichspannung von $U_S = \SI{15}{\volt}$ gewählt.
Das Laborskript sieht hierfür eine Versorgungsspannung zwischen \SI{1}{\volt} und \SI{30}{\volt} vor~\cite{laborskript}.
Bei unveränderten linearen Widerständen bewirkt eine höhere Versorgungsspannung eine größere Diagonalspannung bei gleicher Restverstimmung.
Dadurch lässt sich eine Abweichung vom Nullpunkt bei gleicher absoluter Spannungsauflösung leichter erkennen.
Dem stehen eine höhere Verlustleistung und mögliche temperaturbedingte Widerstandsänderungen gegenüber.
Die Wahl von \SI{15}{\volt} erfolgte als Kompromiss zwischen Abgleichempfindlichkeit und Bauteilbelastung; eine optimale Speisespannung wurde damit nicht nachgewiesen.

\subsection{Bestimmung des Realteils der Impedanz}
Aus den beiden Spannungsteilern der allgemeinen Vergleichsschaltung folgt für den Abgleich $U_d = 0$ die Bedingung
\begin{equation}
  \addxcontentsline{loe}{equation}[\theequation]{Abgleichbedingung der Vergleichsbrücke}
  R_2 R_3 + j \omega R_1 R_2 R_3 C_1
  = R_1 R_4 + j \omega R_1 R_2 R_4 C_2.
  \label{eq:abgleichbedingung}
\end{equation}

Der Vergleich der Realteile von \cref{eq:abgleichbedingung} liefert
\begin{equation}
  \addxcontentsline{loe}{equation}[\theequation]{Berechnung des Parallelwiderstands}
  R_2 = \frac{R_4}{R_3} R_1.
  \label{eq:realgleich}
\end{equation}
Für die AC-Messung mit $\omega \ne 0$ ergibt der Vergleich der Imaginärteile bei von null verschiedenen Widerständen zusätzlich
\begin{equation}
  \addxcontentsline{loe}{equation}[\theequation]{Berechnung der Kapazität bei AC-Abgleich}
  C_2 = \frac{R_3}{R_4} C_1.
  \label{eq:imaggleich}
\end{equation}
Diese zweite Bedingung wird bei der späteren Kapazitätsmessung verwendet; bei Gleichspannung lässt sich $C_2$ aus dem stationären Abgleich nicht bestimmen.

Für die DC-Messung wird $R_1$ so lange verändert, bis die mit einem Multimeter gemessene Diagonalspannung möglichst nahe bei null liegt.
Nach Änderungen der Speisung oder der Widerstände muss der Einschwingvorgang abgewartet werden.
Das Verhältnis von $R_4$ zu $R_3$ wurde so gewählt, dass bei der anschließenden AC-Messung eine mit der vorhandenen Kapazitätsdekade einstellbare Vergleichskapazität verwendet werden kann.
Die eingesetzten Widerstände waren
\begin{itemize}
  \item $R_4 = \SI{470}{\kilo\ohm}$,
  \item $R_3 = \SI{100.1}{\kilo\ohm}$.
\end{itemize}
Daraus ergibt sich das dimensionslose Widerstandsverhältnis
\[
  \frac{R_4}{R_3}
  = \frac{\SI{470}{\kilo\ohm}}{\SI{100.1}{\kilo\ohm}}
  \approx \num{4.695305}.
\]
Für die weiteren Rechnungen wurde das Verhältnis ohne vorzeitiges Runden verwendet.

Da der Einstellbereich der Widerstandsdekade nicht ausreichte, wurde der Vergleichsarm um zwei in Reihe geschaltete Widerstände ergänzt.
Die beiden Zusatzwiderstände wurden auf einem Steckbrett in Reihe geschaltet.
Der gemeinsam gemessene Serienwiderstand beträgt $R_{\mathrm{Serie}} = \SI{1330}{\kilo\ohm}$.

Der gesamte Vergleichswiderstand wird mit $R_1' = R_1 + R_{\mathrm{Serie}}$ bezeichnet.
Bei $R_1' = \SI{1.405}{\mega\ohm}$ wurde eine Restdiagonalspannung von $U_d = \SI{0.01}{\milli\volt}$ gemessen; die Brücke war damit näherungsweise abgeglichen.

In \cref{eq:realgleich} ist für diesen erweiterten Aufbau der gesamte Vergleichswiderstand $R_1'$ einzusetzen.
Damit ergibt sich aus den angegebenen Widerstandswerten
\[
  R_2 = \frac{R_4}{R_3} R_1'
      = \frac{470}{\num{100.1}}\cdot\SI{1.405}{\mega\ohm}
      \approx \SI{6.597}{\mega\ohm}.
\]
Die Rundung dieses Rechenwerts ersetzt keine Bestimmung der Messunsicherheit.

\subsection{Abhängigkeit von der Versorgungsspannung bei nicht abgeglichener Brücke: Plot und Erklärung}
Für die Abhängigkeit von der Versorgungsspannung wurde keine Messreihe aufgenommen.
Ergänzend wurde der Verlauf mit LTspice simuliert und mit MATLAB dargestellt, siehe \cref{fig:dc-bruecke-variable-us}.
Der Plot zeigt somit keine im Labor aufgenommene Messreihe.

\begin{center}
\begin{minipage}{0.95\linewidth}
  \centering
  \includegraphics[width=\linewidth]{dc_bridge_variable_us.png}
  \captionsetup{hypcap=false}
  \captionof{figure}{Abhängigkeit der Differenzspannung von der DC-Versorgungsspannung bei leicht verstimmter Brücke. LTspice-Simulation mit $R_1^{\prime}=\SI{1.405}{\mega\ohm}$, dargestellt mit MATLAB.}
  \label{fig:dc-bruecke-variable-us}
\end{minipage}
\end{center}

Wenn die Widerstände und die Belastung durch das Messgerät gleich bleiben, ist die Differenzspannung $U_d$ proportional zur Versorgungsspannung $U_S$:
\begin{equation}
  \addxcontentsline{loe}{equation}[\theequation]{Abhängigkeit der Differenzspannung von der Versorgungsspannung}
  U_d = k U_S
  \label{eq:bruecke-spannungsabhaengigkeit}
\end{equation}
Der Faktor $k$ hängt von den Widerstandsverhältnissen und dem Eingangswiderstand des Messgeräts ab.
Bei einer nicht abgeglichenen Brücke steigt der Betrag von $U_d$ somit mit der Versorgungsspannung.
Bei ideal abgeglichener Brücke bleibt $U_d=0$.

Der Plot zeigt den erwarteten annähernd linearen Anstieg der Differenzspannung mit der Versorgungsspannung.
In der Simulation wurde der gesamte Vergleichswiderstand mit $R_1'=\SI{1.405}{\mega\ohm}$ eingesetzt.
Bei $U_S=\SI{15}{\volt}$ beträgt die dargestellte Differenzspannung etwa \SI{0.0329}{\milli\volt}; im Labor wurden \SI{0.01}{\milli\volt} gemessen.
Die idealisierte Spannungsteilerrechnung ergibt mit den eingetragenen Widerständen einen Betrag von etwa \SI{0.0319}{\milli\volt}.
Die Simulation liegt damit nahe am idealen Abgleich.
Schon die Rundung des aus dem DC-Abgleich bestimmten Parallelwiderstands auf $R_2=\SI{6.597}{\mega\ohm}$ führt im Modell zu einer kleinen Restspannung.
Die verbleibende Abweichung zum Laborwert kann zusätzlich durch die tatsächlichen Widerstandswerte und die Belastung beziehungsweise den Offset des Messgeräts beeinflusst werden; die einzelnen Beiträge wurden nicht getrennt bestimmt.

\subsection{Auswirkung des Innenwiderstands des Messgeräts}
Die Differenzspannung wurde mit dem Keysight U1232A gemessen.
Das Voltmeter hat einen endlichen Eingangswiderstand. Solange $U_d$ nicht null ist, fließt daher ein kleiner Strom durch das Messgerät.
Dadurch sind die beiden Brückenzweige nicht mehr unabhängig voneinander. Die Berechnung mit zwei unbelasteten Spannungsteilern gilt dann nur näherungsweise.
Die gemessene Differenzspannung fällt durch diese Belastung betragsmäßig kleiner aus als ohne angeschlossenes Messgerät.
Der Eingangswiderstand sollte deshalb deutlich größer als der von der Messdiagonale aus gesehene Innenwiderstand der Brücke sein.

Bei ideal abgeglichener Brücke ist $U_d=0$, und es fließt kein Strom durch das Voltmeter.
Ein rein ohmscher Eingangswiderstand verändert daher die ideale Abgleichbedingung nicht.
Durch die kleinere Anzeige außerhalb des Abgleichs wird es aber schwieriger, eine geringe Verstimmung zu erkennen.



\clearpage
\section{AC-Messbrücke}
Mit der AC-Messbrücke wird die unbekannte Kapazität $C_2$ bestimmt~\cite{laborskript}.

\subsection{Messschaltung}
Die Schaltung aus der DC-Messung wurde übernommen und um die einstellbare Kapazität $C_1$ parallel zu $R_1'$ erweitert, wie in \cref{fig:wien_bridge} dargestellt.
Die Widerstände blieben unverändert: $R_3=\SI{100.1}{\kilo\ohm}$, $R_4=\SI{470}{\kilo\ohm}$ und $R_1'=\SI{1.405}{\mega\ohm}$ einschließlich der Zusatzwiderstände.
Die Brücke wurde mit dem Funktionsgenerator versorgt. Die Differenzspannung wurde mit dem Oszilloskop gemessen.

\subsection{Wahl der Spannung}
Am Funktionsgenerator wurde eine sinusförmige Spannung von $U_{S,\mathrm{pp}}=\SI{4}{\volt}$ eingestellt.
Für einen Sinus ohne Gleichanteil entspricht dies einem Effektivwert von etwa \SI{1.41}{\volt}.
Eine größere Versorgungsspannung macht eine verbleibende Differenzspannung leichter erkennbar, erhöht aber auch die Belastung der Bauteile.

\subsection{Wahl des Frequenzbereichs}
Für die Messung wurde eine Frequenz von $f=\SI{1}{\kilo\hertz}$ eingestellt.
Im Gegensatz zur DC-Messung fließt dabei auch durch die Kondensatoren ein Wechselstrom.
Der kapazitive Blindwiderstand $X_C=-1/(2\pi f C)$ und damit die Empfindlichkeit des Abgleichs hängen von der Frequenz ab.
Mit der einzelnen Messung wurde kein optimaler Frequenzbereich nachgewiesen.

\subsection{Abgleichbedingung der AC-Messbrücke}
Für den idealen Abgleich müssen Real- und Imaginärteil der Differenzspannung null sein.
Aus \cref{eq:abgleichbedingung} folgen für $\omega\neq0$ die beiden Bedingungen
\[
  R_2=\frac{R_4}{R_3}R_1',
  \qquad C_2=\frac{R_3}{R_4}C_1.
\]
Der Widerstandsabgleich wurde aus der DC-Messung übernommen. Zum Kapazitätsabgleich wurde $C_1$ eingestellt.

\subsection{Bestimmung der Kapazität}
Bei einer eingestellten Kapazität von $C_1=\SI{10.6}{\micro\farad}$ wurde am Oszilloskop eine verbleibende Differenzspannung von $U_{\mathrm{diff,rms}}=\SI{10}{\milli\volt}$ gemessen.
Mit \cref{eq:imaggleich} ergibt sich die unbekannte Kapazität zu
\begin{equation}
  \addxcontentsline{loe}{equation}[\theequation]{Berechnung der unbekannten Kapazität}
  C_2=\frac{R_3}{R_4}C_1
     =\frac{\num{100.1}}{470}\cdot\SI{10.6}{\micro\farad}
     \approx\SI{2.26}{\micro\farad}.
  \label{eq:kapazitaet-ergebnis}
\end{equation}
Da die Differenzspannung nicht vollständig null war, ist dies ein Näherungswert aus dem vorgenommenen Abgleich.
Aus der Restspannung allein lässt sich keine genaue Messunsicherheit für $C_2$ angeben.

\clearpage
\section{Subtrahierverstärker}
\subsection{Berechnung der Widerstandswerte}
\subsection{Verstärkungsmessung}
\subsection{Offsetmessung und Eingangs-Offset}
\subsection{CMRR-Messung beziehungsweise Berechnung}
\subsection{Messung der Rauschspannung}

\clearpage
\section{Instrumentenverstärker}
\subsection{Berechnung des Gain-Widerstands}
\subsection{Verstärkungsmessung}
\subsection{Offsetmessung und Eingangs-Offset}
\subsection{CMRR-Messung beziehungsweise Berechnung}
\subsection{Vergleich von Subtrahier- und Instrumentenverstärker}

\clearpage
\section{Gewichtsmessung mit Halbbrücke}
\subsection{Berechnung der Sensitivität (Aufgabe 3.4)}
\subsection{Messung der Sensitivität und Plot}
\subsection{Offsetmessung}
\subsection{Messung des Handy-Gewichts}
\subsection{Vergleich mit der Waage}

\clearpage
\section{Temperatursensitivität der Halbbrücke}
\subsection{Messung der RMS-Rauschspannung}
\subsection{Plot der Verstärkung über der Temperatur}
\subsection{Plot der Rauschspannung über der Temperatur}

\clearpage
\section{Vollbrücke}
\subsection{Sensitivitätsmessung und Vergleich mit der Halbbrücke}
\subsection{Temperaturabhängigkeit}
\subsection{Plot der Rauschspannung über der Temperatur}

\clearpage
\section{Tiefpassfilter zur Rauschunterdrückung}
\subsection{Messung der gefilterten Rauschspannung}
\subsection{Erklärung der Rauschquellen}

\clearpage
% Weitere Quellen als \bibitem{schluessel} ergänzen und mit \cite{schluessel} zitieren.
\renewcommand{\refname}{Literaturverzeichnis}
\begin{thebibliography}{99}
  \bibitem{vorlage-laborprotokoll}
  \texttt{bendemaj}: \emph{Laborprotokoll -- LaTeX-Vorlage für Laborprotokolle}.
  GitHub-Repository.
  \url{https://github.com/bendemaj/Laborprotokoll}.
  Zugriff am 14.~November~2026.

  \bibitem{laborskript}
  \emph{Messtechnik Labor: Skriptum}, Wintersemester 2026/2027.
  Abschnitt 3.1.1, S.~29--30 (PDF-S.~33--34).

  % Hier später weitere Literaturangaben einfügen.
\end{thebibliography}
\end{document}
